% Public text-only snapshot of the instructor's Module 12 GAN slide deck.
% Upstream: /Users/setup/Library/CloudStorage/Box-Box/Teaching/6050/LaTeX/Module 12/12.4S GAN.tex
% Upstream SHA-256: 705c63ba8e34f0954d85614fc08701b3a86b4fcf698be29b89e1f4c127b65a8f
% Local preamble dependencies, TikZ source, executable code, external assets,
% prescriptive hyperparameters, product-current comparisons, and claims without a
% sufficiently clear verification boundary are omitted. The book independently checks
% and derives every published claim, implementation, experiment, and figure.
\documentclass{article}
\usepackage{amsmath}
\usepackage{amssymb}
\usepackage[margin=1in]{geometry}
\usepackage{hyperref}

\title{Lecture 12.4: Generative AI and Generative Adversarial Networks}
\author{Heman Shakeri}
\date{}

\begin{document}
\maketitle

\section*{Public Snapshot Boundary}

This standalone, text-only snapshot preserves the instructor-authored GAN sequence and
selected equations. Diagrams, executable examples, method-ranking claims, quotations,
and practical recipes have been removed. The retained notes document provenance and
must be checked against primary sources before publication.

\section{The Adversarial Setup}

A generator $G$ maps noise $\mathbf{z}\sim p_{\mathbf{z}}$ to a sample
$G(\mathbf{z})$. A discriminator $D$ maps an observation to a scalar interpreted,
under the original formulation, as the probability that the observation came from the
data rather than the generator. Training couples the two models: the discriminator
learns from real and generated observations, while the generator learns through the
discriminator.

The lecture names Goodfellow et al.'s ``Generative Adversarial Nets'' as its primary
starting point: \url{https://arxiv.org/abs/1406.2661}.

% [The original adversarial-game and network diagrams are omitted.]

\section{The Minimax Objective}

The original two-player value function is

\[
\min_G\max_D V(D,G)
=\mathbb{E}_{\mathbf{x}\sim p_{\mathrm{data}}}
[\log D(\mathbf{x})]
+\mathbb{E}_{\mathbf{z}\sim p_{\mathbf{z}}}
[\log(1-D(G(\mathbf{z})))].
\]

For a fixed generator, the discriminator is trained to distinguish observations drawn
from $p_{\mathrm{data}}$ from those drawn from the generator distribution $p_g$. In
practice, the generator is often trained with the non-saturating objective

\[
\max_G\;
\mathbb{E}_{\mathbf{z}\sim p_{\mathbf{z}}}
[\log D(G(\mathbf{z}))],
\]

which has the same target equilibrium as the original minimax generator objective but
different optimization behavior.

\section{Alternating Optimization and the Idealized Equilibrium}

A conceptual training iteration alternates two updates:

\begin{enumerate}
\item hold $G$ fixed and update $D$ using real and generated observations;
\item hold $D$ fixed and update $G$ through the current discriminator.
\end{enumerate}

Under the idealized analysis of the original objective, an optimal discriminator for a
fixed generator is

\[
D_G^*(\mathbf{x})
=\frac{p_{\mathrm{data}}(\mathbf{x})}
{p_{\mathrm{data}}(\mathbf{x})+p_g(\mathbf{x})}.
\]

If $p_g=p_{\mathrm{data}}$, this expression gives
$D_G^*(\mathbf{x})=1/2$ wherever the common density is supported. This theoretical
statement assumes the relevant densities and optimization idealizations; it is not a
diagnostic that finite neural-network training has converged.

% [The original fixed update-count recommendations are omitted; useful schedules
% depend on the models, optimizer, data, and diagnostics.]

\section{Architectural Pattern}

For image data, the lecture describes a generator that progressively increases spatial
resolution and a discriminator that progressively reduces it. The output layer and
data scaling must agree, and normalization choices differ between the two networks.
These are architectural patterns rather than universal requirements.

% [Generator and discriminator diagrams, dimension tables, and implementation-specific
% prescriptions are omitted.]

\section{Why Training Can Fail}

The lecture uses three failure modes to organize GAN diagnostics:

\begin{itemize}
\item \emph{mode collapse}: different latent inputs map to a limited set of outputs;
\item \emph{saturation or weak gradients}: the current discriminator supplies an
unhelpful learning signal to the generator;
\item \emph{oscillation or imbalance}: alternating updates chase a moving opponent
rather than settling near a useful solution.
\end{itemize}

The generator's gradient depends on the discriminator's current geometry, so stronger
classification performance is not automatically a better generator-learning signal.
The original deck's simplified causal account of vanishing gradients is omitted; the
book must distinguish the minimax and non-saturating losses and rederive their gradients
before making that argument.

% [The original instability and mode-collapse illustrations are omitted.]

\section{Historical Extensions}

The lecture names convolutional GAN design, Wasserstein objectives, gradient penalties,
conditional GANs, and unpaired image translation as historically important extensions.
These approaches change different parts of the contract: architecture, divergence or
distance surrogate, regularization, conditioning information, or supervision.

For unpaired translation between domains $X$ and $Y$, a cycle-consistency term can be
written as

\[
\mathcal{L}_{\mathrm{cyc}}(G,F)
=\mathbb{E}_{x\sim p_X}
[\|F(G(x))-x\|_1]
+\mathbb{E}_{y\sim p_Y}
[\|G(F(y))-y\|_1].
\]

This term augments adversarial objectives in both domains. It encodes a consistency
preference; it does not by itself identify a uniquely correct correspondence between
unpaired observations. The displayed formulation follows Zhu et al.,
\href{https://openaccess.thecvf.com/content_iccv_2017/html/Zhu_Unpaired_Image-To-Image_Translation_ICCV_2017_paper.html}
{``Unpaired Image-to-Image Translation Using Cycle-Consistent Adversarial
Networks''}.

\section{A Comparative Role}

GANs, variational autoencoders, and diffusion models use different training signals and
offer different forms of inference or evaluation. A method comparison is meaningful
only after fixing the data, task, compute budget, metrics, and implementation. The
lecture's release-era ranking and categorical account of why one family ``won'' are
therefore omitted.

% [Product-current method comparisons, unverified quotations, named-product examples,
% and practical hyperparameter recipes are omitted.]

\section{Lasting Concepts}

The adversarial formulation established a reusable idea: learn a generator through a
learned critic rather than through a fixed pointwise reconstruction loss alone. It also
made failure analysis central to generative modeling---coverage, optimization dynamics,
and evaluation cannot be inferred from sample appearance alone.

\end{document}
